\begin{ouverture}[college2]{11}{Poids, masse et gravitation}{Partie 2 · Mouvements et interactions}{Pourquoi pèse-t-on moins sur la Lune ?}
\objectif{Distinguer masse et poids, instruments et unités.}
\objectif{Utiliser une proportionnalité entre poids et masse.}
\objectif{Décrire qualitativement la gravitation et ses limites de modèle.}
\prerequis{Force, masse et proportionnalité.}
\end{ouverture}
\pcPage{Le cours}{Poids, masse et gravitation}
\phantomsection\label{ess-11}
\essSection{1}{Une distinction indispensable}
La masse $m$, mesurée avec une balance et exprimée en kilogrammes, ne dépend pas du lieu dans les situations du collège. Le poids $P$ est la force de pesanteur exercée sur l’objet ; il se mesure avec un dynamomètre et s’exprime en newtons. Près de la surface terrestre, il est vertical et dirigé vers le centre de la Terre dans le modèle sphérique.\par
\pcFigure{figures/c11-terrelune.pdf}{Même objet : masse identique, poids différent selon l’astre. Les flèches représentent les poids, pas les masses.}
\essSection{2}{Une relation de proportionnalité}
En un lieu donné, le poids est proportionnel à la masse : $P=mg$. Le coefficient $g$ s’exprime en N/kg. On prend souvent $g\approx10\unite{N/kg}$ sur Terre pour les calculs simples, ou $9{,}8\unite{N/kg}$ si l’énoncé le précise ; sur la Lune, environ $1{,}6\unite{N/kg}$. Convertir la masse en kilogrammes avant de calculer. Ces valeurs sont des approximations locales.\par
\clearpage
\essSection{3}{Mesurer et exploiter un graphique}
Suspendre plusieurs masses connues permet de tracer $P$ en fonction de $m$. Une droite passant près de l’origine traduit la proportionnalité ; le quotient $P/m$ donne une estimation de $g$. Un écart de mesure se discute avant de déclarer la loi fausse. La droite doit être construite à partir des mesures, sans relier mécaniquement tous les points par une ligne brisée.\par
\pcFigure{figures/c11-graphique.pdf}{Données idéales pour comparer les pentes sur Terre et sur la Lune ; une vraie mesure présenterait une dispersion.}
\essSection{4}{La gravitation et la suite du cycle}
La gravitation est une interaction attractive à distance entre objets possédant une masse. Elle intervient dans les mouvements des planètes et des satellites. Plus les masses sont grandes, plus l’action est forte ; l’éloignement la diminue. Dans ce manuel, on travaille surtout qualitativement. L’exploitation numérique de la loi de gravitation fournie et la structure de l’atome seront approfondies en 3e. L’impesanteur apparente d’astronautes en orbite ne signifie pas que la Terre n’agit plus sur eux.\par
\begin{pcEncadre}{method}{Un exemple expliqué}Un sac a une masse de $600\unite g=0{,}600\unite{kg}$. Avec $g=10\unite{N/kg}$, son poids terrestre est $P=0{,}600\times10=6{,}0\unite N$. Sur la Lune, $P=0{,}600\times1{,}6=0{,}96\unite N$. Le sac ne perd pas de matière pendant ce changement de lieu : sa masse reste $0{,}600\unite{kg}$.\end{pcEncadre}
\pcPage{Expérimenter et retenir}{Poids, masse et gravitation}
\phantomsection\label{activite-11}
\begin{pcEncadre}{method}{Investigation en classe}\textbf{Question.} Poids et masse sont-ils proportionnels ? Utiliser des masses scolaires, une balance et un dynamomètre adapté. Relever au moins quatre couples $(m;P)$, en notant m en kg et P en N. Tracer le graphique et estimer $P/m$ pour chaque couple. Repérer un décalage du zéro éventuel. \textbf{À rendre :} tableau, droite d’ajustement visuelle, estimation de $g$ et limites. Le déplacement vers la Lune est étudié par données documentaires, pas par une manipulation.\end{pcEncadre}
\pcSousTitre{Ce que je dois savoir expliquer}\begin{questions}
\item La masse s’exprime en kg et le poids en N.
\item Le poids est une force, avec direction et sens.
\item À lieu donné, $P=mg$.
\item Le coefficient $g$ dépend du lieu.
\item La gravitation agit à distance et reste présente en orbite.
\end{questions}
\begin{pcEncadre}{attention}{Une erreur à éviter}Ne pas utiliser $m=600$ dans $P=mg$ si la masse est donnée en grammes et $g$ en N/kg : il faut d’abord écrire $m=0{,}600\unite{kg}$.\end{pcEncadre}
\pcRepere{cor-activite-11}{Pistes pour l’investigation}\pcGrille
\clearpage\markright{Je démarre}
\rubrique{Je démarre}
\par\noindent\begin{minipage}{\linewidth}
\begin{exercice}[1]{Associer les unités}
\begin{questions}
\item Associer kilogramme et newton à masse et poids.
\item Quel instrument mesure le poids ?
\end{questions}
\end{exercice}
\end{minipage}
\par\noindent\begin{minipage}{\linewidth}
\begin{exercice}[1]{Le lieu change}
\begin{questions}
\item La masse d’un objet change-t-elle lors d’un voyage Terre-Lune dans le modèle du cours ?
\item Et son poids ?
\end{questions}
\end{exercice}
\end{minipage}
\par\noindent\begin{minipage}{\linewidth}
\begin{exercice}[1]{Convertir une masse}
\begin{questions}
\item Convertir $250\unite g$ en kg.
\item Convertir $1{,}5\unite{kg}$ en g.
\end{questions}
\end{exercice}
\end{minipage}
\par\noindent\begin{minipage}{\linewidth}
\begin{exercice}[1]{Un poids simple}
\begin{questions}
\item Calculer le poids de $2\unite{kg}$ avec $g=10\unite{N/kg}$.
\item Donner sa direction et son sens près du sol.
\end{questions}
\end{exercice}
\end{minipage}
\par\noindent\begin{minipage}{\linewidth}
\begin{exercice}[1]{Une interaction}
\begin{questions}
\item La gravitation est-elle une action de contact ?
\item Est-elle attractive dans les situations étudiées ?
\end{questions}
\end{exercice}
\end{minipage}
\clearpage\markright{Je m’entraîne}
\rubrique{Je m’entraîne}
\par\noindent\begin{minipage}{\linewidth}
\begin{exercice}[2]{Un petit objet}
\begin{questions}
\item Calculer le poids terrestre d’un objet de $80\unite g$ avec $g=10\unite{N/kg}$.
\item Quelle erreur ferait-on en oubliant la conversion ?
\end{questions}
\end{exercice}
\end{minipage}
\par\noindent\begin{minipage}{\linewidth}
\begin{exercice}[2]{Retrouver la masse}
\begin{questions}
\item Un objet a un poids de $35\unite N$ avec $g=10\unite{N/kg}$. Quelle masse ?
\item Exprimer en grammes.
\end{questions}
\end{exercice}
\end{minipage}
\par\noindent\begin{minipage}{\linewidth}
\begin{exercice}[2]{Estimer le coefficient}
\begin{questions}
\item Une masse de $0{,}50\unite{kg}$ pèse $4{,}9\unite N$. Estimer $g$.
\item Quelle unité utiliser ?
\end{questions}
\end{exercice}
\end{minipage}
\par\noindent\begin{minipage}{\linewidth}
\begin{exercice}[2]{Terre et Lune}
\begin{questions}
\item Une masse de $5\unite{kg}$ est étudiée avec $g_T=10$ et $g_L=1{,}6\unite{N/kg}$. Calculer les poids.
\item Quel rapport entre les poids ?
\end{questions}
\end{exercice}
\end{minipage}
\par\noindent\begin{minipage}{\linewidth}
\begin{exercice}[2]{Une droite}
\pcFigureEx{figures/ex-poids-graphe.pdf}{}
\begin{questions}
\item Les couples $(m;P)$ sont $(0{,}1;1)$, $(0{,}2;2)$ et $(0{,}3;3)$, en kg et N. Quel modèle ?
\item Prévoir le poids de $0{,}45\unite{kg}$.
\end{questions}
\end{exercice}
\end{minipage}
\clearpage\markright{Je raisonne}
\rubrique{Je raisonne}
\par\noindent\begin{minipage}{\linewidth}
\begin{exercice}[3]{Une balance détournée}
\begin{questions}
\item Un pèse-personne terrestre transforme une force mesurée en une indication en kg. Pourquoi cette indication peut-elle être fausse sur la Lune ?
\item Quel appareil comparerait des masses sans cette conversion locale ?
\end{questions}
\end{exercice}
\end{minipage}
\par\noindent\begin{minipage}{\linewidth}
\begin{exercice}[3]{La phrase du quotidien}
\begin{questions}
\item Une personne dit « je pèse 50 kg ». Reformuler scientifiquement.
\item Quel serait son poids avec $g=10\unite{N/kg}$ ?
\end{questions}
\end{exercice}
\end{minipage}
\par\noindent\begin{minipage}{\linewidth}
\begin{exercice}[3]{Un zéro décalé}
\pcFigureEx{figures/ex-poids-decale.pdf}{Exemple de relevé expérimental présentant un décalage de zéro.}
\begin{questions}
\item Un graphique expérimental de P en fonction de m coupe l’axe vertical à 0,3 N. Quelle cause possible ?
\item Doit-on modifier la loi immédiatement ?
\end{questions}
\end{exercice}
\end{minipage}
\par\noindent\begin{minipage}{\linewidth}
\begin{exercice}[3]{Une orbite}
\begin{questions}
\item Des astronautes flottent dans une station en orbite. Peut-on conclure que la gravitation a disparu ?
\item Quelle idée qualitative explique leur situation ?
\end{questions}
\end{exercice}
\end{minipage}
\par\noindent\begin{minipage}{\linewidth}
\begin{exercice}[3]{Une donnée de gravitation}
\begin{questions}
\item Un modèle documentaire indique que doubler une distance entre deux objets divise la force gravitationnelle par quatre. Si elle valait 12 N, quelle nouvelle valeur ?
\item Cette règle donne-t-elle à elle seule le poids de tout objet ?
\end{questions}
\end{exercice}
\end{minipage}
\pcPage{Défi du chapitre}{Poids, masse et gravitation}
\phantomsection\label{defi-ex-11}
\pcDefiTitre{11}{Le colis lunaire}
\begin{questions}
\item Un colis a une masse de $750\unite g$. Calculer son poids terrestre avec $g_T=9{,}8\unite{N/kg}$ et lunaire avec $g_L=1{,}6\unite{N/kg}$.
\item Choisir un calibre parmi 1 N, 5 N et 10 N pour la mesure terrestre.
\item Si le dynamomètre affiche 0 N en chute libre, que peut-on conclure sur la masse ?
\end{questions}
\pcGrille
